\FrontChapter{References}\label{references}
The references below identify classical definitions, proof inputs and adjacent
constructions. They do not constitute an exhaustive priority review. The
arithmetic seed and the Gamma--Laguerre comparisons inherited from the
project are documented in the preserved research archive.

\begingroup
\makeatletter
\patchcmd{\thebibliography}{\chapter*{\bibname}}{}{}{\PackageError{RC}{Bibliography heading patch failed}{}}
\makeatother
\begin{thebibliography}{15}
\setlength{\itemsep}{11pt}
\setlength{\parskip}{3pt}
\bibitem{Mendez2024} Miguel A. Méndez.
\emph{Combinatorial differential operators in: Faà di Bruno formula,
enumeration of ballot paths, enriched rooted trees and increasing rooted trees}.
Australasian Journal of Combinatorics 89(3) (2024), 350--384.
Example 2.4 contains the diagonal feedback equation.
\url{https://ajc.maths.uq.edu.au/pdf/89/ajc_v89_p350.pdf}

\bibitem{Dascaliuc2025} Radu Dascaliuc, Tuan N. Pham, Enrique Thomann and Edward C. Waymire.
\emph{Transformation of Stochastic Recursions and Critical Phenomena in the
Analysis of the Aldous-Shields-Athreya Cascade and Related Mean Flow Equations}.
arXiv:2411.00629v2 (2025).
The nonlinear equation $u'=-u+u(\alpha t)^2$ is an adjacent model, not the
equation studied here.
\url{https://arxiv.org/abs/2411.00629v2}

\bibitem{Coffey2006} Mark W. Coffey.
\emph{On the coefficients of the Baez-Duarte criterion for the Riemann
hypothesis and their extensions}. arXiv:math-ph/0608050 (2006).
Its signed Möbius weights differ from the positive squarefree weights used here.
\url{https://arxiv.org/abs/math-ph/0608050}

\bibitem{DLMF} NIST Digital Library of Mathematical Functions.
Sections 5.5 and 5.12 give the gamma reflection and beta identities.
Section 10.46 defines the Mittag--Leffler function. Sections 16.8 and 16.21
give the generalized hypergeometric and Meijer $G$ differential equations.
Section 18.2 develops general orthogonal-polynomial theory.
\url{https://dlmf.nist.gov/}

\bibitem{Uvarov1969} V. B. Uvarov.
\emph{The connection between systems of polynomials that are orthogonal
with respect to different distribution functions}.
USSR Computational Mathematics and Mathematical Physics 9(6) (1969), 25--36.
\url{https://doi.org/10.1016/0041-5553(69)90124-4}

\bibitem{Elkies} Noam D. Elkies.
\emph{Functions of finite order: product formula and logarithmic derivative}.
Harvard Math 259 course notes. The finite-order Hadamard theorem supplies the
classical factorization input.
\url{https://people.math.harvard.edu/~elkies/M259.02/prod.pdf}

\bibitem{Byrnes2013} Alyssa Byrnes, Lin Jiu, Victor H. Moll and Christophe Vignat.
\emph{Recursion Rules for the Hypergeometric Zeta Functions}.
arXiv:1305.1892 (2013). Hypergeometric zero spectra and coefficient recursions
provide related constructions.
\url{https://arxiv.org/abs/1305.1892}

\bibitem{Voros2003} André Voros.
\emph{Zeta functions for the Riemann zeros}.
Annales de l'Institut Fourier 53(3) (2003), 665--699.
Sections 2.1--2.2 develop related canonical-product and Mellin methods.
\url{https://doi.org/10.5802/aif.1955}

\bibitem{Voros2004} André Voros.
\emph{Erratum: Zeta functions for the Riemann zeros}.
Annales de l'Institut Fourier 54(4) (2004), 1139.
\url{https://doi.org/10.5802/aif.2046}

\bibitem{OEIS} The OEIS Foundation Inc.
\emph{The On-Line Encyclopedia of Integer Sequences}, sequence A000182.
This is the tangent-number anchor. No unmatched-prefix novelty claim or
submission to OEIS is asserted.
\url{https://oeis.org/A000182}

\bibitem{Alrabaiah2020} Hussam Alrabaiah, Israr Ahmad, Kamal Shah and Ghaus Ur Rahman.
\emph{Qualitative analysis of nonlinear coupled pantograph differential
equations of fractional order with integral boundary conditions}.
Boundary Value Problems 2020, article 138.
This is a coupled fractional boundary-value system with different hypotheses.
\url{https://doi.org/10.1186/s13661-020-01432-2}

\bibitem{Blasiak2011} Pawel Blasiak and Philippe Flajolet.
\emph{Combinatorial Models of Creation--Annihilation}.
Séminaire Lotharingien de Combinatoire 65 (2011), B65c.
Proposition 24 and Note 25 discuss multitype increasing trees and differential systems.
\url{https://arxiv.org/abs/1010.0354v3}

\bibitem{Mauldin1988} R. Daniel Mauldin and S. C. Williams.
\emph{Hausdorff Dimension in Graph Directed Constructions}.
Transactions of the American Mathematical Society 309(2) (1988), 811--829.
Its spectral-radius balance describes a geometric dimension. The analytic
order law here requires its own proof.
\url{https://sites.math.unt.edu/~mauldin/papers/no67.pdf}

\bibitem{Knessl2006} Charles Knessl and Wojciech Szpankowski.
\emph{Binary Trees, Left and Right Paths, WKB Expansions, and Painlevé Transcendents}.
Proceedings of ANALCO (2006), 198--204.
The ordinary generating-function identity has different normalization and
coefficient rules from the differential system here.
\url{https://www.cs.purdue.edu/homes/spa/papers/analco06.pdf}

\bibitem{Schilling2010} René L. Schilling, Renming Song and Zoran Vondraček.
\emph{Bernstein Functions: Theory and Applications}.
First edition. De Gruyter, 2010. ISBN 978-3-11-021530-4.
Proposition 1.2 and Theorem 1.4 give the classical Laplace uniqueness and
Bernstein representation results.
\url{https://doi.org/10.1515/9783110215311}
\end{thebibliography}
\endgroup
